\section{Immutable Parent Pointers}
\label{sec:immutable-parent-pointers}

The immutable parent pointer is the foundational structural primitive of the Recursive Spawning Protocol. It is the single mechanism that transforms the delegation chain from a forensic reconstruction attempted after the fact into a runtime-verifiable artifact --- one that is immutable, auditable, and architecturally enforced. It renders all modes of autonomous drift structurally impossible regardless of operational urgency, architectural complexity, or adversarial sophistication of any machine actor in the chain.

This section specifies the parent pointer formally (\S\ref{sec:parent-pointer-formal}), defines the chain traversal operator and establishes the base case for the root human anchor (\S\ref{sec:chain-definition}), specifies immutability enforcement through the Fact Layer write protocol (\S\ref{sec:fact-layer-enforcement}), describes the Purpose Layer's spawn authorization role (\S\ref{sec:purpose-layer-authorization}), and provides the complete chain traversal algorithm with correctness arguments (\S\ref{sec:chain-traversal}).

% ------------------------------------------------------------------------- %
\subsection{Formal Definition: ParentPointer}
\label{sec:parent-pointer-formal}

\begin{dfn}[Parent Pointer]\label{dfn:parent-pointer}
Let $\mathcal{N}$ be the set of all agent nodes in a \bxthree{} system implementing the Recursive Spawning Protocol, and let $P$ be the set of all Purpose Layer actors. The \emph{parent pointer} $\alpha(n)$ is defined for every node $n \in \mathcal{N}$ as the total function
\begin{equation}
\alpha : \mathcal{N} \rightarrow P \cup \{\bot\}
\label{eq:alpha-function}
\end{equation}
where $\alpha(n) = p$ if $p \in P$ is the Purpose Layer actor that authorized the spawn event resulting in $n$'s activation, and $\alpha(n) = \bot$ if and only if $n$ is the root human anchor established at system initialization. The pointer is established exactly once at node provisioning via a Fact Layer atomic write operation and satisfies the following immutability property:
\begin{equation}
\forall n,\ p:\ \bigl(\alpha(n) = p \ \land \ \textit{provisioned}(n)\bigr) \implies \neg\exists\ \textit{modify}(\alpha(n))
\label{eq:parent-pointer-immutable}
\end{equation}
No actor --- including the Purpose Layer after provisioning, the Bounds Engine, any administrative process, or any external principal --- may execute an operation that overwrites, redirects, or deletes $\alpha(n)$ after the provisioning write is confirmed. The Fact Layer write protocol rejects any such attempt at the protocol level, before execution reaches the ledger.
\end{dfn}

The parent pointer is not a reference variable held by the parent node toward its child. It is a Fact Layer property intrinsic to the child node, established at provisioning time and stored in the node's sealed memory envelope. The directionality is architecturally critical: the child carries $\alpha(n)$ as an immutable attribute of itself, not as a mutable external reference held by the parent. This means the chain is traversable from any node outward to its complete lineage regardless of the operational state of any other node in the chain. A parent that has terminated, crashed, or been decommissioned does not affect the child's ability to reference its parent through $\alpha(n)$.

The parent pointer is distinct from a conventional access-controlled pointer in a second critical respect. Access control enforces immutability by restricting which principals may issue modification operations; if sufficient privilege is acquired through credential theft, software vulnerability, or administrative compromise, the access control barrier is circumvented. The Fact Layer write protocol enforces immutability by making the modification operation structurally impossible --- there is no defined mechanism for processing a \textit{modify}($\alpha(n)$) request regardless of any privilege held. There is no privileged override, no emergency bypass, and no administrative role with the capability to alter $\alpha(n)$ after provisioning. This is not a stronger access control policy; it is a different architectural category.

The initialization of $\alpha(n)$ occurs atomically during the spawn event. The Fact Layer records the spawn event in the forensic ledger and writes $\alpha(n) = p$ into the node's sealed memory envelope as part of the same atomic transaction. This separation of authorization (Purpose Layer) from assignment (Fact Layer) is essential: it prevents a parent from claiming to have spawned a node whose pointer was set to a different actor, and it prevents a child from falsifying its own ancestry after provisioning.

\begin{dfn}[Spawn Event]\label{dfn:spawn-event}
A \emph{spawn event} is the 5-tuple
\begin{equation}
\mathrm{spawn}(c,\ p,\ \sigma_c,\ t,\ \phi)
\label{eq:spawn-event}
\end{equation}
where $c$ is the newly created child node, $p \in P$ is the authorizing Purpose Layer actor, $\sigma_c$ is the authorized execution scope, $t$ is the wall-clock timestamp, and $\phi$ is the cryptographic signature produced by $p$ over the event payload using $p$'s Purpose Layer private key. The signature $\phi$ attests that $p$ has authorized the spawn and binds $c$ to $p$ in the forensic ledger. The Fact Layer verifies $\phi$ before accepting the spawn event for ledger recording.
\end{dfn}

The spawn event is the sole legitimate mechanism for assigning a parent pointer. There is no other path by which $\alpha(c)$ can be set. This eliminates any ambiguity about the origin of a node's parent pointer: if $\alpha(c) = p$ appears in the ledger, it exists because a spawn event bearing $p$'s cryptographic signature was recorded, and no other origin is possible.

% ------------------------------------------------------------------------- %
\subsection{Chain Definition and Base Case}
\label{sec:chain-definition}

\begin{dfn}[Chain Operator]\label{dfn:chain-operator}
For any node $a \in \mathcal{N}$, the \emph{chain operator} $\Chain{a}$ is defined recursively as
\begin{equation}
\Chain{a} \equiv
\begin{cases}
\emptyset & \text{if }\ \alpha(a) = \bot \quad \text{(base case: root human)} \\[6pt]
\{\, \ParentPointer(p,\ a) \, \} \cup \Chain{p} & \text{if }\ \alpha(a) = p \neq \bot \quad \text{(recursive case)}
\end{cases}
\label{eq:chain-operator}
\end{equation}
where $\ParentPointer(p,\ a)$ denotes the ordered pair $(p,\ a)$ establishing that $p$ is the immutable parent of $a$. The empty set $\emptyset$ returned in the base case indicates that the root human has no parent pointer --- it is the architectural terminus of all chains.
\end{dfn}

The chain operator produces an ordered set of parent-child pairs representing the complete ancestral lineage of $a$, terminating at the root human. The recursion terminates because $\alpha$ is a total function and each recursive step applies $\alpha$ to the parent node, which is strictly closer to the root by depth. The depth of any node $n$ is defined as:
\begin{equation}
\mathrm{depth}(n) \equiv
\begin{cases}
0 & \text{if }\ \alpha(n) = \bot \\
1 + \mathrm{depth}(\alpha(n)) & \text{otherwise}
\end{cases}
\label{eq:depth-definition}
\end{equation}

The depth function is well-defined because $\alpha$ has no cycles --- the Fact Layer rejects any spawn event that would introduce a cycle, and $\alpha$ is total. Every node has exactly one parent unless it is the root human, and the root human has null parent.

\begin{dfn}[Root Human Anchor]\label{dfn:root-human}
The root human anchor $h \in P$ is the unique Purpose Layer actor designated at system initialization as the terminus of all delegation chains. It satisfies $\alpha(h) = \bot$ and $\mathrm{depth}(h) = 0$. No machine actor can serve as a root: only a human principal designated at initialization holds this role.
\end{dfn}

The root human anchor is non-collapsing and non-substitutable. As chains grow in depth through successive spawn events, the root remains fixed. When a child node $n_i$ spawns a grandchild $n_{i+1}$, the root for $n_{i+1}$ is the same root as for $n_i$: the root does not advance to intermediate machine actors. This is what makes the upstream accountability guarantee structurally operative regardless of chain depth.

The uniqueness of the chain for any node follows directly from the uniqueness of $\alpha(n)$ for each $n$. Because $\alpha$ is a total function with a unique value for each node, $\Chain{a}$ is uniquely determined for any input $a$. There is no ambiguity, no alternative lineage, and no path through which $a$ could be traced to a different root.

% ------------------------------------------------------------------------- %
\subsection{Immutability Enforcement: The Fact Layer Write Protocol}
\label{sec:fact-layer-enforcement}

The immutability of $\alpha(n)$ is enforced by the Fact Layer write protocol, which defines exactly two write operations relevant to parent pointers:

\begin{enumerate}
    \item[\textbf{W1}] $\mathrm{WriteParentPointer}(c,\ p)$: The initial provisioning write, executed exactly once during a spawn event, recording $\alpha(c) = p$.
    \item[\textbf{W2}] No subsequent write operation targeting $\alpha(n)$ for any $n$ where $\textit{provisioned}(n) = \text{true}$.
\end{enumerate}

The Fact Layer write protocol processes write requests in the following order: (1) authenticate the requestor's identity, (2) verify the request type against the defined write operation set, (3) execute the write if the type is W1 and all preconditions are satisfied, or (4) reject with error code \texttt{E-NOT-PERMITTED} if the type is W2 or any other undefined operation targeting $\alpha(n)$. Steps (1) and (2) are applied to every write request without exception; there is no bypass path, no debug mode, and no emergency override that disables step (2).

The enforcement is cryptographic as well as protocol-level. After W1 is executed, the parent pointer $\alpha(c)$ and associated metadata are encrypted inside the node's sealed memory envelope using a Fact Layer-only symmetric key. The envelope is then HMAC-signed with a Fact Layer-only key. Any subsequent read of $\alpha(c)$ requires Fact Layer decryption; any modification attempt invalidates the HMAC and is rejected at signature verification. The sealed memory envelope provides defense-in-depth: even if a software vulnerability allowed a malformed write request to bypass the protocol-level rejection, the envelope's cryptographic integrity check would reject the modification.

The sealed memory envelope mechanism operates as follows:

\begin{enumerate}
    \item \textbf{Encrypted at rest.} $\alpha(n)$ is encrypted in the node's persistent memory envelope using a Fact Layer-only symmetric key. Even a compromised Bounds Engine cannot read $\alpha(n)$ because it never has access to the decryption key.
    \item \textbf{HMAC-signed.} The envelope carries an HMAC-SHA-256 signature computed over its contents using a Fact Layer-only key. Any modification to the ciphertext invalidates the signature.
    \item \textbf{Decrypted only for reads.} The Fact Layer decrypts the envelope only when reading $\alpha(n)$ to service a chain-traversal or audit request. The plaintext is never exposed to the Bounds Engine or the child node's own code.
    \item \textbf{Re-sealed after reads.} After each authorized read, the Fact Layer re-encrypts and re-signs the envelope, preventing replay attacks based on stale plaintext.
\end{enumerate}

The Fact Layer operates as a standalone, isolated process with its own memory space, inaccessible to the child node's address space. The combination of the protocol-level write constraint and the sealed envelope ensures that $\alpha(n)$ is modification-resistant in all deployment contexts.

\medskip
\noindent\textbf{Why access control is insufficient.} The distinction between access-control-based immutability and protocol-level immutability is the distinction between a promise and a guarantee. In an access-control model, immutability is enforced by restricting which principals may issue modification operations; if an attacker acquires sufficient privileges through credential theft, insider compromise, or software vulnerability, the barrier is circumvented and immutability is violated. The immutability was a policy, enforceable until it was not.

In the Fact Layer write protocol, the immutability of $\alpha(n)$ is a property of the protocol itself. There is no operation defined for modifying $\alpha(n)$ after provisioning, so no amount of privilege elevation, credential compromise, or software vulnerability can produce a valid modification operation. The protocol does not enforce immutability --- it makes immutability the only valid state transition. The distinction is architectural, not a matter of degree.

This is the precise failure mode the Fact Layer write protocol eliminates: retroactive manipulation of chain topology for the purpose of accountability laundering. Without protocol-level immutability, an adversarial actor could re-parent a node to a favorable parent after the fact, obscuring the actual delegation chain and insulating the true authorizing principal from scrutiny. With the Fact Layer write protocol, this attack is structurally impossible because no valid operation exists to produce it.

% ------------------------------------------------------------------------- %
\subsection{Spawn Authorization: The Purpose Layer}
\label{sec:purpose-layer-authorization}

The Purpose Layer serves as the exclusive authorization origin for all parent pointer assignments. It evaluates every spawn request against three mandatory conditions before issuing a warrant, and it never delegates this evaluation to any subordinate machine actor.

\begin{dfn}[Purpose Layer Spawn Authorization]\label{dfn:purpose-authorization}
A parent node $n_i \in \mathcal{N}$ that intends to spawn a child $n_{i+1}$ must submit a spawn request to the Purpose Layer declaring: (P1) the proposed child operating scope $R(n_{i+1})$, (P2) the specific operational condition necessitating spawning, and (P3) a demonstration that the condition cannot be resolved within $R(n_i)$. The Purpose Layer validates three mandatory conditions before issuing a spawn warrant:

\begin{enumerate}
    \item[(P1)] \textbf{Scope refinement.} $R(n_{i+1})$ must be a strict subset of $R(n_i)$:
    \begin{equation}
    R(n_{i+1}) \subset R(n_i).
    \label{eq:scope-refinement}
    \end{equation}
    Lateral scope stability --- where $R(n_{i+1}) = R(n_i)$ --- is not authorized. The chain grows through scope \emph{narrowing}, not sideways expansion. Equality at any link constitutes a drift condition and invalidates the warrant.

    \item[(P2)] \textbf{Operational necessity.} The parent must declare, in the Fact Layer authorization ledger, the precise condition requiring child spawning. The declaration must identify why the condition cannot be resolved within $R(n_i)$ and why the child is the minimum necessary decomposition of the current task.

    \item[(P3)] \textbf{Minimum scope proportionality.} The child scope must be the narrowest scope sufficient to address the declared condition. The Purpose Layer does not authorize child scopes broader than operationally necessary.
\end{enumerate}
\end{dfn}

When all three conditions are satisfied, the Purpose Layer issues a spawn warrant $w_i$ --- a cryptographically signed record written atomically to the Fact Layer authorization ledger:
\begin{equation}
w_i = \bigl\langle
\text{parent}=n_i,\;
\text{child}=n_{i+1},\;
R(n_{i+1}),\;
\text{justification},\;
\text{timestamp},\;
\sigma_{\text{PL}}
\bigr\rangle .
\label{eq:spawn-warrant}
\end{equation}

The warrant is the sole authorized precondition for the Fact Layer's WriteParentPointer operation. No machine actor may initiate a spawn event without a matching warrant in the Fact Layer ledger --- the Fact Layer pre-commit hook rejects any provisioning request that lacks a valid warrant before the operation reaches the ledger write stage. This makes the Purpose Layer's exclusive terminus role a structural guarantee, not a configuration option.

The Purpose Layer also enforces two structural constraints at spawn authorization:

\begin{enumerate}
    \item[(C1)] \textbf{Cycle prevention.} The proposed parent $n_i$ must not be a descendant of the proposed child $n_{i+1}$ in the current chain. This is verified by checking $n_{i+1} \notin \Chain{n_i}$. If the proposed child already appears in the parent's ancestry chain, the spawn is denied because it would introduce a cycle in $\alpha$, violating the acyclicity of the delegation graph.
    \item[(C2)] \textbf{Depth bound compliance.} The depth of the proposed child must satisfy $\mathrm{depth}(n_{i+1}) \leq D_{\max}$, where $D_{\max}$ is the architectural maximum depth parameter recorded in the Fact Layer authorization ledger at system initialization. If the spawn would cause the depth bound to be exceeded, the spawn is refused and the node is placed in \texttt{ESCALATING} state with a mandatory Bailout Protocol trigger.
\end{enumerate}

The Purpose Layer's exclusive authority over spawn authorization means that no machine actor in the chain can autonomously extend the chain. Every link in the delegation chain is Purpose Layer-authorized before it is Fact Layer-recorded. This two-layer gate --- Purpose Layer warrant first, Fact Layer recording second --- is the fundamental control structure that makes the RSP's accountability guarantees operative.

% ------------------------------------------------------------------------- %
\subsection{Chain Traversal Algorithm}
\label{sec:chain-traversal}

The chain traversal algorithm provides the runtime mechanism for computing $\Chain{a}$ from any node $a$, and the chain verification algorithm provides the runtime procedure for confirming structural validity of the chain to a human principal.

\begin{algorithm}
\caption{Chain-Traverse($n$) --- Build the delegation chain from node $n$ to the human anchor}
\label{alg:chain-traverse}
\begin{algorithmic}[1]
\REQUIRE $n$ is a provisioned RSP node with $\alpha$ defined
\ENSURE $\Chain{n}$ --- the ordered list of all ancestor nodes including the human anchor
\STATE $chain \leftarrow \langle \rangle$ \COMMENT{Initialize empty ordered list}
\STATE $current \leftarrow n$ \COMMENT{Start at the query node}
\WHILE{$\mathrm{true}$}
    \STATE $\mathrm{append}(chain,\ current)$ \COMMENT{Add current node to chain}
    \IF{$\alpha(current) = \bot$}
        \STATE $\textbf{break}$ \COMMENT{Root human anchor reached --- chain complete}
    \ENDIF
    \STATE $current \leftarrow \alpha(current)$ \COMMENT{Move to parent}
\ENDWHILE
\RETURN $chain$
\end{algorithmic}
\end{algorithm}

\begin{algorithm}
\caption{Chain-Verify($n$) --- Verify the authorization chain from node $n$ to human anchor}
\label{alg:chain-verify}
\begin{algorithmic}[1]
\REQUIRE $n$ is a provisioned RSP node
\ENSURE $\mathrm{Boolean}$ indicating whether the chain is a structurally valid RSP chain
\STATE $chain \leftarrow \mathrm{Chain\text{-}Traverse}(n)$
\STATE $m \leftarrow \mathrm{len}(chain)$

\STATE \COMMENT{\textit{Step V1: Verify chain terminates at human anchor}}
\IF{$\alpha(chain[m-1]) \neq \bot$}
    \STATE \RETURN $\text{false}$ \COMMENT{Root is not $\bot$ --- malformed chain}
\ENDIF
\IF{$\neg h(chain[m-1])$}
    \STATE \RETURN $\text{false}$ \COMMENT{Root is not a designated human principal}
\ENDIF

\STATE \COMMENT{\textit{Step V2: Verify each spawn event has a Purpose Layer warrant}}
\FOR{$i \leftarrow 0$ \TO $m-2$}
    \STATE $child \leftarrow chain[i]$
    \STATE $parent \leftarrow chain[i+1]$
    \IF{$\neg \text{warrant\_exists}(child,\ parent)$}
        \STATE \RETURN $\text{false}$ \COMMENT{Missing Purpose Layer warrant for spawn}
    \ENDIF
\ENDFOR

\STATE \COMMENT{\textit{Step V3: Verify scope refinement at each link}}
\FOR{$i \leftarrow 0$ \TO $m-2$}
    \STATE $child \leftarrow chain[i]$
    \STATE $parent \leftarrow chain[i+1]$
    \IF{$\neg (R(child) \subset R(parent))$}
        \STATE \RETURN $\text{false}$ \COMMENT{Scope refinement constraint violated}
    \ENDIF
\ENDFOR

\STATE \COMMENT{\textit{Step V4: Verify depth bound compliance at each node}}
\FOR{$i \leftarrow 0$ \TO $m-1$}
    \IF{$\mathrm{depth}(chain[i]) > D_{\max}$}
        \STATE \RETURN $\text{false}$ \COMMENT{Depth bound exceeded at node}
    \ENDIF
\ENDFOR

\RETURN $\text{true}$
\end{algorithmic}
\end{algorithm}

Chain-Verify returns \texttt{true} if and only if all four verification conditions hold simultaneously. V1 confirms the chain terminates at a human anchor; V2 confirms every spawn was authorized by the Purpose Layer; V3 confirms scope refinement was respected at every link; V4 confirms the depth bound was respected throughout. A chain that passes Chain-Verify is, by the RSP definitions, a fully authorized delegation chain terminating at a human principal.

The algorithm is deterministic and produces the same result regardless of which node initiates verification. There is no self-serving interpretation of authorization that a drifted or orphaned node could produce, because Chain-Verify inspects the complete chain topology and warrants in the Fact Layer ledger, not the node's own claims about its authorization.

\medskip
\noindent\textbf{Correctness arguments.} The correctness of Chain-Traverse follows from three structural properties: (1) $\alpha$ is a total function, so $\alpha(current)$ is defined for every provisioned node; (2) the loop invariant guarantees that at each iteration, $current$ is the next node in the chain toward the root; (3) the loop terminates when $\alpha(current) = \bot$, which occurs at the root human by Definition~\ref{dfn:root-human}, and the depth bound $D_{\max}$ guarantees this occurs within $D_{\max} + 1$ iterations. The output is uniquely determined because $\alpha$ has a unique value for each node.

The correctness of Chain-Verify follows from the completeness of its four verification conditions: V1 checks the termination condition (base case), V2 checks every link for Purpose Layer authorization, V3 checks every link for scope refinement compliance, and V4 checks every node for depth bound compliance. If any condition fails, the chain does not satisfy the RSP definition and the function returns \texttt{false}. If all conditions pass, the chain satisfies all RSP structural invariants and the function returns \texttt{true}.

\medskip
\noindent\textbf{Constant-time escalation bound.} The depth bound $D_{\max}$ provides a constant-time upper bound on chain traversal cost. The worst-case number of pointer dereferences required to reach the human anchor from any node is $D_{\max} + 1$, independent of the total number of nodes in the system. A system with ten thousand spawned nodes and $D_{\max} = 5$ requires at most six pointer dereferences to reach a human --- the same as a system with ten nodes. Mutable parent pointers would allow the chain to be extended retroactively, making the effective depth potentially unbounded and the worst-case escalation time unbounded as well. Immutability is what makes the depth bound a structural guarantee, not merely a policy.

In the Bailout Protocol context, chain traversal is used to identify the correct Human Accountability Anchor when an agent at any depth encounters an unresolvable exception. The upward propagation of the exception signal follows the same parent-pointer structure as chain traversal, but uses the asynchronous trigger pathway rather than the synchronous verification pathway. Chain traversal ensures that this walk never diverges, never shortcuts, and never fails to reach a human.

\medskip
\noindent\textbf{Computational cost.} Chain-Traverse runs in $O(d)$ time where $d$ is the chain depth (at most $D_{\max} + 1$). Chain-Verify runs in $O(m)$ for the traversal plus $O(m)$ for each of the three verification loops, yielding $O(m)$ total where $m$ is chain length, also bounded by $D_{\max} + 1$. Both algorithms are linear in chain length and do not depend on the total number of nodes in the system.

% ------------------------------------------------------------------------- %
\subsection{Summary: Immutability as Architectural Constraint}
\label{sec:immutability-summary}

The immutable parent pointer is the load-bearing constraint of the Recursive Spawning Protocol. Without it, the chain is mutable --- subject to retroactive modification by actors with sufficient privilege --- and the RSP's correctness properties (drift prevention, chain integrity, termination) collapse to policy conventions rather than structural guarantees.

The immutability is architectural --- not a stronger access control policy --- because it is enforced by the Fact Layer write protocol, which defines no valid modification operation for $\alpha(n)$ after provisioning. This is distinct from access-control-based immutability in four ways: there is no override path for any actor; the immutability holds under all system failure conditions including those that disable other enforcement mechanisms; all sources receive equivalent enforcement treatment regardless of privilege level; and the parent pointer and its Purpose Layer warrant are created atomically with no temporal window for inconsistency.

The chain operator $\Chain{a}$ is uniquely determined by the immutable parent pointers in the chain. The chain traversal algorithm computes this operator correctly and terminates. Chain-Verify provides the runtime verification procedure for chain structural validity. Together, these mechanisms make the delegation chain an immutable, verifiable, and auditable runtime artifact --- the foundation on which all other RSP properties depend.